\section{在项目中使用前端技术}
\subsection{学习内容对应具体任务}
我在本次实践中的具体方向是站点、内容与文件，包括公共导航、首页与品牌展示、玩法文章、FAQ、下载中心，%
以及内容和媒体管理页面\cite{local-team}。这些页面既服务于阅读内容的访客，也服务于维护内容的管理员。%


在公共布局中，我接触到组件复用、路由和响应式样式；在文章页面中，我需要理解接口读取、%
异步状态和图文排版；在管理界面中，我又进一步接触表单、文件、选区和版本冲突。%
表\ref{tab:used}整理了当前工程实际使用的技术及其位置。

\begin{table}[!htb]
\centering
\smalltable
\caption{本项目实际使用的前端技术}%
\label{tab:used}
\begin{tabularx}{\linewidth}{>{\raggedright\arraybackslash}p{2.75cm}YY}
\toprule
技术 & 具体用途 & 能力边界\\
\midrule
Vue 3 与组合式 API & 公共导航、文章列表、编辑器和状态组件 & 组织界面与逻辑；不能代替后端数据校验\\
\tabledivider
TypeScript & 内容、媒体、文件等接口类型；组件参数 & 提供编译期约束；外部响应仍需运行时判断\\
\tabledivider
Vue Router 与 Pinia & 页面导航、路由懒加载、会话共享 & 管理端入口判断不等于服务端授权\\
\tabledivider
Element Plus & 管理页面中的现成对话框等界面组件 & 提供组件库能力，与 Vue 的组件机制不是同一层\\
\tabledivider
CSS Grid / Flexbox & 图文并列、卡片网格、窄屏单列 & 布局正确仍需在具体视口下检查\\
\tabledivider
Fetch 与异步逻辑 & 读取文章、筛选资料、提交内容 & 请求可能乱序、失败或结果暂时不确定\\
\tabledivider
Three.js CSS3DRenderer & 拍立得商品卡片的倾斜与回弹 & 变换的是 DOM 元素，不是实体商品三维模型\\
\tabledivider
浏览器动画与观察 API & 页面入场、可见性判断、逐帧运动 & 特性或设备条件不足时保留静态内容\\
\tabledivider
原生编辑能力与自制区块 & 文字格式、左右图文、图注和预览 & 支持受控内容结构，未实现多人实时协作\\
\bottomrule
\end{tabularx}
\end{table}

\subsection{三个让我继续深入调研的问题}

第一个问题来自商品展示：一张图片怎样看起来像可以轻轻拿起的相纸？%
最开始我把这种立体感理解为“制作一个三维模型”，但代码中的 CSS3DRenderer
变换的其实是包含图片和文字的 DOM。这个差别使我继续查阅 CSS 三维、WebGL 和 WebGPU，%
理解不同层次的“网页三维”。

第二个问题来自内容页面：打开文章时，页面代码、文章数据和图片为什么不是同时出现？%
项目的路由使用动态 import，详情页又通过接口读取文章，正文图片还设置了懒加载。%
这说明“加载页面”至少包含加载代码、读取数据、获取图片和完成绘制几个过程。%


第三个问题来自后台编辑：为什么看起来只是一块可以输入文字的区域，却要处理图注、粘贴、%
段落结构和重复保存？当前编辑器既要让管理员方便操作，又要保证前台能稳定展示。%
因此，我从 contenteditable 延伸到结构化编辑器，进一步了解内容模型、事务和协作机制。%
这三个问题构成后续调研的主线。

\subsection{从页面效果理解组件和状态}
在 Vue 中，我逐渐形成“界面是当前状态的结果”的理解，可简写为
\[
\text{界面}=F(\text{数据},\text{交互状态}).
\]
例如，同一个文章区域可以显示加载提示、成功内容、空结果或错误提示。组件化让这些状态拥有统一表现，%
但并不会自动替我决定失败后保留哪些输入、是否允许重试。项目中的 ContentState、ContentBody 和
ContentEditor 分别承担状态展示、正文渲染与编辑职责；把它们区分开后，%
排查问题会比让一个页面处理所有细节清楚得多。
